\subsection{Data Sources}
\label{data:sources}
The data come from three published studies of the human gut microbiome, covering colorectal cancer (CRC)
\citep{wirbel2019}, inflammatory bowel disease (IBD) \citep{lloydprice2019} and graft-versus-host disease
(GvHD) \citep{peled2020}. Their clinical endpoints play no role here. What each study contributes is a set
of bacterial genera, a set of viral operational taxonomic units (vOTUs), and abundance and protein
profiles for both, measured in the same samples.

The CRC data derive from \citet{wirbel2019}, a meta-analysis of eight faecal metagenomic studies of colorectal
cancer conducted in Germany, China, the United States and Japan, covering 768 cases together with controls. Of
these, 84 samples carry paired bacterial and viral measurements and enter this thesis.

The IBD data derive from \citet{lloydprice2019}, a longitudinal study of 132 subjects with Crohn's disease,
with ulcerative colitis, or without inflammatory bowel disease, each followed for one year at up to 24 time
points. Of the specimens collected there, 148 samples enter this thesis.

The GvHD data derive from \citet{peled2020}, a study of 1,362 patients undergoing allogeneic
haematopoietic-cell transplantation at four centres in New York, North Carolina, Germany and Japan, from which
8,767 faecal samples were collected. Of these, 90 samples enter this thesis.

Each cohort provides four tables: abundances for the bacterial genera and for the vOTUs, a spacer table
indexed by genus and by vOTU, and the protein profiles of both. We report their sizes with the other
descriptive quantities in the results.

Write \(\mathcal{B}=\{b_1,\ldots,b_n\}\) for the bacterial genera and
\(\mathcal{V}=\{v_1,\ldots,v_m\}\) for the vOTUs of a cohort that carry both a protein profile and an
abundance profile. Both networks below are indexed by the same candidate pairs, the Cartesian product
\begin{equation}\label{eq:pairs}
  \mathcal{P}=\mathcal{B}\times\mathcal{V},
\end{equation}
so a pair \((b_i,v_j)\) is one element of \(\mathcal{P}\) regardless of which network is being described.

\subsection{Protein Clusters}
\label{sec:mmseqs}

Individual protein sequences do not give a common vocabulary across organisms, so we first group similar
sequences into protein clusters with MMseqs2 \citep{steinegger2017}. MMseqs2 prefilters candidate
sequence pairs, aligns the kept candidates, and clusters those that are sufficiently similar; the same
prefilter-then-align strategy is used by other large-scale protein aligners \citep{buchfink2021}.

The clusters give one protein-count matrix per node set,
\[
  X^{b}\in\mathbb{R}^{n\times d_b},
  \qquad
  X^{v}\in\mathbb{R}^{m\times d_v},
\]
whose rows are organisms and whose columns are clusters; entry \((i,k)\) is the number of proteins of
organism \(i\) assigned to cluster \(k\). The two vocabularies are built separately and are of very
different size, so they do not share coordinates.

\subsection{Pair Feature}
\label{sec:phi}

The transformation \(\varphi\) of Section~\ref{sec:pipeline} maps the profiles \(x^{b}_{i}\) and
\(x^{v}_{j}\) of the two nodes of a pair to a single vector \(z_{ij}\in\mathbb{R}^{d}\). We fix it once,
before fitting any model, and use the same transformation for every model and both networks.

We use the protein clusters that occur in both vocabularies,
\(S=S^{\mathrm{all}}_{b}\cap S^{\mathrm{all}}_{v}\), so that coordinate \(k\) is the same cluster for the
bacterium and for the vOTU. We reduce the counts to presence,
\(\operatorname{bin}(x)_k=\mathbf{1}\{x_k>0\}\), which is invariant to per-organism rescaling and
therefore unaffected by proteome size or assembly completeness, and add the two halves coordinatewise,
\begin{equation}\label{eq:phi}
  z_{ij}=\operatorname{bin}\!\bigl(x^{b}_{i,S}\bigr)+\operatorname{bin}\!\bigl(x^{v}_{j,S}\bigr)
  \;\in\;\{0,1,2\}^{d}.
\end{equation}
So \(z_{ij,k}\) counts how many of the two partners carry cluster \(k\): none, one, or both. 

\subsection{CRISPR Network}
\label{sec:crispr}
If the spacers of a bacterium are compared against the genomes of viruses, a sufficiently strong match provides 
evidence that the bacterium, or a related bacterial lineage, has previously encountered the corresponding virus 
or a closely related virus. The absence of a match, however, does not establish the opposite.

One of the main issues is that CRISPR spacers are only around 30 nucleotides long, so standard nucleotide 
searches can lose sensitivity, especially when the viral sequence has diverged from the sequence from which the 
spacer was originally acquired. SpacePHARER \citep{zhang2021spacepharer} addresses this by translating both 
CRISPR spacers and viral genomes in all six reading frames and comparing the resulting sequences at the protein 
level. Candidate hits are subsequently realigned at the nucleotide level. SpacePHARER then combines the evidence 
from multiple spacers of one bacterium against the same viral genome into a single score and controls the 
resulting matches using a false discovery rate threshold. The search was run separately for each cohort, and 
its output enters as the CRISPR table of Section~\ref{data:sources}.

Only bacteria carrying a CRISPR--Cas system produce spacers, the occurrence of such systems differs between
taxonomic groups, and spacers are lost over evolutionary time. An interaction leaving no spacer record is
therefore indistinguishable from one that never occurred, and \(Y\) records sequence-supported historical
relationships rather than all interactions.


The network is defined over all of \(\mathcal{P}\), with \(Y_{ij}=1\) when the search reports CRISPR
evidence for the pair \((b_i,v_j)\) and \(Y_{ij}=0\) otherwise.

\subsection{Abundance Network}

Marginal association between two taxa can be induced by a third, so the network is defined by conditional
association instead. For variables \(a=(a_1,\ldots,a_P)\) with covariance \(\Sigma\), the precision matrix
\(\Theta=\Sigma^{-1}\) satisfies \(\Theta_{k\ell}=0\) exactly when \(a_k\) and \(a_\ell\) are
conditionally independent given all remaining variables, under a Gaussian assumption \citep{Bishop2006}.
Its non-zero off-diagonal entries are the edges of the network.

SPIEC-EASI \citep{kurtz2015} estimates such a network from compositional abundance data in three steps.
The abundance matrix \(C\in\mathbb{R}_{\geq0}^{N\times P}\) of \(N\) samples and \(P\) variables is first
centred-log-ratio transformed \citep{aitchison1982},
\(\tilde c_{sk}=\log c_{sk}-P^{-1}\sum_{\ell}\log c_{s\ell}\). Its precision matrix is then estimated
under an \(L_{1}\) penalty by the graphical lasso \citep{friedman2008},
\begin{equation}\label{eq:glasso}
  \widehat\Theta_{\lambda}
  =\operatorname*{arg\,min}_{\Theta\succ0}
  \Bigl\{\operatorname{tr}\bigl(\widehat\Sigma\Theta\bigr)-\log\det\Theta
  +\lambda\textstyle\sum_{k\neq\ell}\lvert\Theta_{k\ell}\rvert\Bigr\},
\end{equation}
where a larger \(\lambda\) gives a sparser network. Finally \(\lambda\) is chosen by StARS
\citep{liu2010}: the estimate is refitted on \(B\) subsamples, the selection frequency of an edge is
\begin{equation}\label{eq:stars-pi}
  \widehat\pi^{\,\mathrm{gl}}_{k\ell}(\lambda)
  =\frac{1}{B}\sum_{r=1}^{B}
  \mathbf{1}\bigl\{\widehat\Theta^{(r)}_{\lambda,k\ell}\neq0\bigr\},
\end{equation}
and \(\widehat\lambda_{\mathrm{StARS}}\) is the least regularised value at which the average edge
instability \(2\widehat\pi^{\,\mathrm{gl}}(1-\widehat\pi^{\,\mathrm{gl}})\) stays below a threshold
\(\xi^{\ast}\). We apply SPIEC-EASI to each cohort separately with \(\xi^{\ast}=0.02\) rather than the
standard \(0.05\), since the selected network is unstable at less conservative thresholds.

The abundance table stacks both domains, so \(P=n+m\) and \(\widehat\Theta\) is one symmetric
\((n+m)\times(n+m)\) matrix with three blocks,
\begin{equation}\label{eq:blocks}
  \widehat\Theta=
  \begin{pmatrix}
    \widehat\Theta^{\,\mathcal{B}\mathcal{B}} & \widehat\Theta^{\,\mathcal{B}\mathcal{V}}\\[2pt]
    \bigl(\widehat\Theta^{\,\mathcal{B}\mathcal{V}}\bigr)^{\!\top} & \widehat\Theta^{\,\mathcal{V}\mathcal{V}}
  \end{pmatrix},
  \qquad
  \widehat\Theta^{\,\mathcal{B}\mathcal{V}}\in\mathbb{R}^{n\times m},
\end{equation}
of which we keep only the off-diagonal block and discard the within-domain ones. The estimated network is
therefore not bipartite but a rectangular sub-block of a conditional-dependence network over all \(n+m\)
taxa, and a kept entry describes a conditional dependence given every other taxon, including those whose
blocks are dropped.

The weighted network is not an entry of \(\widehat\Theta\) but the StARS selection frequency of the
corresponding cross-block edge at the selected penalty. Writing \(k(i)\) and \(\ell(j)\) for the indices
of \(b_i\) and \(v_j\) among the \(P\) variables,
\begin{equation}\label{eq:Wdef}
  W_{ij}
  =\widehat\pi^{\,\mathrm{gl}}_{k(i)\,\ell(j)}\bigl(\widehat\lambda_{\mathrm{StARS}}\bigr)\in[0,1],
\end{equation}
so \(W_{ij}\) is the proportion of subsamples in which that edge was selected: a measure of how
reproducibly the edge is recovered, not a partial correlation or a precision entry.

\subsection{Missing Values and \(W\) Binarisation}
\label{sec:mask}

The two networks are constructed independently and do not span the same nodes. The CRISPR network is
defined over all of \(\mathcal{P}\), whereas the abundance network covers only the taxa that entered
SPIEC-EASI, so \(W\) is missing for the remaining pairs. With
\(\mathcal{B}_W\subseteq\mathcal{B}\) and \(\mathcal{V}_W\subseteq\mathcal{V}\) the nodes it contains,
\begin{equation}\label{eq:maskdef}
  M_{ij}=\mathbf{1}\{b_i\in\mathcal{B}_W\}\,\mathbf{1}\{v_j\in\mathcal{V}_W\},
  \qquad
  \mathcal{O}=\bigl\{(i,j)\in\mathcal{P}:M_{ij}=1\bigr\},
\end{equation}
and pairs outside \(\mathcal{O}\) are missing values of \(W\) rather than absent edges. The indicator
factorises into a row term and a column term, so whether a pair is observed is determined by its two
nodes alone and not by the value of \(W\). We fit and evaluate every model on \(\mathcal{O}\), a complete-case
analysis, so the design matrix has \(|\mathcal{O}|\) rows rather than \(nm\).

\paragraph{\(W\) Binarisation}
For binary prediction, thresholding at \(\tau\in[0,1)\) gives
\begin{equation}\label{eq:thresh}
  W^{(\tau)}_{ij}=\mathbf{1}\{W_{ij}>\tau\},
\end{equation}
so \(W^{(0)}_{ij}\) marks every pair selected in at least one subsample and a larger \(\tau\) demands a
more reproducibly selected edge. We use \(\tau=0\) in all later analyses. This is the natural cut between
an edge that is never selected and one that is selected at all, it requires no tuning, and it leaves a
near-balanced split between the two classes, whereas a larger \(\tau\) leaves too few positive pairs to
fit. In continuous form \(W_{ij}\) is bounded while the regression models predict on the real line, so we
fit on a logit scale and map back before scoring: we clip values to \([\varepsilon,1-\varepsilon]\) with
\(\varepsilon=10^{-3}\), transform by \(\widetilde W_{ij}=\operatorname{logit}(\cdot)\), and return the
fitted values as \(\widehat W_{ij}=\sigma(\widehat{\widetilde W}_{ij})\), against which we compute
\(R^{2}\) on the original scale.

\subsection{Model Hyperparameters}

The model and optimization hyperparameters are summarized in Table~\ref{tab:model_hyperparameters}. Both neural models are
regularized with dropout \citep{srivastava2014} and trained with the AdamW optimizer \citep{loshchilov2019}.

\begin{table}[htbp]
\centering
\caption{Hyperparameters used for the baseline and latent-space models.}
\label{tab:model_hyperparameters}
\begin{tabular}{lll}
\hline
\textbf{Model} & \textbf{Hyperparameter} & \textbf{Value} \\
\hline

Sparse logistic regression
& penalty & \(L_1\) \\
& regularization & selected by 3-fold CV \\
& selection rule & one-standard-error rule \\
& solver & liblinear \\
& maximum iterations & 2000 \\
& class weighting & balanced \\
\hline

Ridge regression
& \(L_2\) penalty \(\alpha\) & 1.0 \\
\hline

XGBoost classifier
& number of trees & 500 \\
& maximum tree depth & 4 \\
& learning rate & 0.05 \\
& row subsampling & 0.8 \\
& feature subsampling & 0.8 \\
& class weighting & \(n_{\mathrm{neg}}/n_{\mathrm{pos}}\) \\
\hline

XGBoost regression
& number of trees & 500 \\
& maximum tree depth & 4 \\
& learning rate & 0.05 \\
& row subsampling & 0.8 \\
& feature subsampling & 0.8 \\
\hline

MLP classifier/regressor
& hidden dimension & 64 \\
& activation & ReLU \\
& dropout & 0.3 \\
& batch size & 512 \\
& optimizer & AdamW \\
& learning rate & \(3\times10^{-4}\) \\
& weight decay & \(10^{-3}\) \\
& epochs & 200 \\
\hline

Joint latent model
& joint latent dimension & 64 \\
& head weights & \(\gamma_Y=\gamma_W=1\) \\
& activation & ReLU \\
& dropout & 0.3 \\
& batch size & 512 \\
& optimizer & AdamW \\
& learning rate & \(3\times10^{-4}\) \\
& weight decay & \(10^{-3}\) \\
& epochs & 200 \\
\hline
\end{tabular}
\end{table}
\FloatBarrier
